% =====================================================================
%  Lecture 02 -- 2G Circuit Switching                     (English edition)
% =====================================================================
\documentclass[11pt,a4paper]{article}
\providecommand\ACNroot{../..}
\input{\ACNroot/templates/lecture-en}

\ACNset{number}{02}
\ACNset{title}{2G Circuit Switching}
\ACNset{shorttitle}{2G Circuit Switching}
\ACNset{titleother}{سوئیچ مداری در نسل دوم}
\ACNset{part}{Part I -- 2G/3G foundations (Lectures 1--4)}
\ACNset{time}{6--8 hours (3 h reading, 3--5 h exercises)}
\ACNset{level}{Intermediate}
\ACNset{prereq}{Lecture 1 (GSM architecture); basic probability (Poisson process, exponential distribution)}
\ACNset{gen}{2}

\begin{document}
\ACNcover

% ---------------------------------------------------------------------
\section{Learning objectives}
After studying this lecture you should be able to:
\begin{objectives}
  \item \textbf{Explain} what a circuit is, which resources a GSM call reserves and why
        reservation guarantees quality but wastes capacity.
  \item \textbf{Describe} the SS7 protocol stack (MTP, SCCP, TCAP, MAP, ISUP, BSSAP) and
        the role of signalling points and signal transfer points.
  \item \textbf{Trace} a mobile-originated and a mobile-terminated call, including HLR
        interrogation and the use of the MSRN.
  \item \textbf{Derive} the Erlang-B formula from the $M/M/m/m$ model and
        \textbf{dimension} a cell or trunk group for a target blocking probability.
  \item \textbf{Quantify} trunking efficiency and the cost of splitting a channel pool.
  \item \textbf{Analyse} roaming signalling and call routing, and \textbf{evaluate} why
        circuit switching became a limitation for data.
\end{objectives}

% ---------------------------------------------------------------------
\section{Prerequisite knowledge}
You need the GSM architecture and identities of Lecture~1 (MSC, VLR, HLR, GMSC, IMSI,
TMSI) and elementary probability: a Poisson arrival process with rate $\lambda$, and an
exponentially distributed holding time with mean $h$. We use the word \emph{traffic} in
the telephony sense: the average number of simultaneously busy circuits.

\paragraph{Key terms} The vocabulary of this lecture with its Persian equivalents:
\begin{keyterms}
Circuit switching & \fa{سوئیچ مداری} & A dedicated path is reserved end to end for the whole call. \\
Signalling System No.~7 (SS7) & \fa{سامانهٔ سیگنالینگ شمارهٔ ۷} & Out-of-band signalling network between exchanges. \\
ISUP & \fa{بخش کاربر ISDN} & SS7 protocol for setting up and releasing trunk circuits. \\
MAP & \fa{بخش کاربرد موبایل} & SS7 application protocol for mobility (HLR, VLR, MSC). \\
MSISDN & \fa{شمارهٔ تلفن مشترک} & The dialled telephone number of a subscriber. \\
MSRN & \fa{شمارهٔ رومینگ ایستگاه موبایل} & Temporary number used to route a call to the serving MSC. \\
Erlang (E) & \fa{ارلانگ} & Unit of traffic: one circuit busy continuously. \\
Blocking probability & \fa{احتمال انسداد} & Probability that a call attempt finds all circuits busy. \\
Grade of service & \fa{درجهٔ خدمت} & Target blocking probability used for dimensioning. \\
Trunking efficiency & \fa{بازده ترانکینگ} & Carried traffic per circuit; grows with pool size. \\
Roaming & \fa{رومینگ} & Use of a visited network with a home subscription. \\
\end{keyterms}

% ---------------------------------------------------------------------
\section{Introduction and motivation}
Telephony was built on one simple promise: once the called party answers, the voice
path behaves like a pair of wires, with constant delay, no loss and no competition from
other users. The engineering way to keep that promise is to \emph{reserve} resources for
the whole call. GSM inherited this model from ISDN, because its first job was to be a
mobile extension of the fixed telephone network and because 1990s exchanges, trunks and
billing systems were all built around 64\,kbit/s circuits.

Two problems had to be solved on top of the classic telephone network. First,
\emph{signalling} had to carry much more than digits: authentication, location updates,
handover, roaming between countries. Second, the network had to \emph{find} a subscriber
whose location changes, before any circuit can be set up. GSM solved both with the
Signalling System No.~7 (SS7) and its Mobile Application Part (MAP). This lecture follows
a call from dial to hang-up, then asks the quantitative question every operator faces:
\emph{how many circuits are enough?}

% ---------------------------------------------------------------------
\section{Fundamental concepts}

\subsection{What a circuit reserves}
For a GSM voice call the network reserves, for its entire duration:
\begin{itemize}
  \item one traffic channel (TCH) on the radio interface (a timeslot on a carrier);
  \item one 16\,kbit/s sub-channel on the Abis link and one 64\,kbit/s timeslot on the A
        interface after transcoding;
  \item a connection through the switching matrix of the MSC;
  \item a 64\,kbit/s trunk on each inter-exchange link (one ISUP circuit per link).
\end{itemize}
The reservation is held even during silence. In a typical conversation each direction
is active only about 40\% of the time, so most of the reserved capacity carries nothing.

\begin{definitionbox}{Traffic intensity}
If calls arrive at rate $\lambda$ (calls per second) and each lasts on average $h$ seconds,
the offered traffic is $A = \lambda h$, a dimensionless number measured in erlangs. One
erlang is the traffic that keeps one circuit busy all the time.
\end{definitionbox}

\subsection{In-band versus common-channel signalling}
Early telephone networks signalled with tones on the voice circuit itself. This was slow,
limited in vocabulary and insecure (tones could be imitated by users). Common-channel
signalling separates the two: signalling messages travel over a dedicated packet network,
while voice circuits carry only voice. SS7 is the ITU-T standard for this signalling
network \cite{ituq700}. It is the first example in this course of the
\emph{separation of control plane and user plane}, an idea that will return in every
generation.

\subsection{Signalling points}
Nodes of an SS7 network are \emph{signalling points} (SPs), each identified by a point
code. A \emph{service switching point} (SSP) is an exchange such as an MSC. A
\emph{signal transfer point} (STP) is a packet router of the SS7 network: it forwards
messages between SPs without terminating them, and STPs are deployed in mated pairs for
reliability. A \emph{service control point} (SCP) hosts databases and service logic,
for example the intelligent-network logic used for prepaid charging.

% ---------------------------------------------------------------------
\section{Detailed technical explanation}

\subsection{The SS7 protocol stack}
\Cref{fig:ss7} shows the stack as used in GSM.
\begin{itemize}
  \item \textbf{MTP 1--3} (Message Transfer Part) carries messages reliably between
        adjacent nodes (MTP\,2: framing, error detection and retransmission) and routes them
        across the signalling network using point codes (MTP\,3).
  \item \textbf{SCCP} (Signalling Connection Control Part) adds connection-less and
        connection-oriented services and, crucially for mobility, \emph{global title
        translation}: a message can be addressed with a telephone-number-like global title
        (for example derived from an IMSI) and the network finds the point code of the HLR.
  \item \textbf{TCAP} (Transaction Capabilities Application Part) structures
        request--response dialogues between applications.
  \item \textbf{MAP} \cite{ts29002} runs over TCAP and defines mobility operations:
        \msg{Update Location}, \msg{Insert Subscriber Data}, \msg{Send Authentication Info},
        \msg{Send Routing Information} (SRI), \msg{Provide Roaming Number} (PRN),
        \msg{Cancel Location}.
  \item \textbf{ISUP} (ISDN User Part) runs directly over MTP\,3 and sets up and releases
        trunk circuits between exchanges with messages such as IAM, ACM, ANM, REL and RLC.
  \item \textbf{BSSAP} runs over SCCP on the A interface. Its BSSMAP part carries messages
        between MSC and BSC (assignment, handover, ciphering); its DTAP part transparently
        carries messages between MSC and MS (call control and mobility management).
\end{itemize}

\begin{acnfigure}
  \acnfig{lecture-02/ss7-stack}
  \caption{SS7 stack in a GSM core network. MTP provides transport; ISUP controls trunk
  circuits; SCCP and TCAP support the MAP mobility dialogues and BSSAP on the A interface.}
  \label{fig:ss7}
\end{acnfigure}

\subsection{Layer-3 protocols between MS and MSC}
On the radio side, layer 3 is divided into three sublayers \cite{ts24008}: \emph{radio
resource management} (RR, terminated in the BSS: channel assignment, handover),
\emph{mobility management} (MM: location updating, authentication, TMSI reallocation) and
\emph{connection management} (CM, with call control CC: Setup, Alerting, Connect,
Disconnect). MM and CC messages are carried transparently through the BSS in DTAP.

\subsection{Numbers and identities used for routing}
Three numbers must not be confused:
\begin{itemize}
  \item \textbf{MSISDN} -- the number people dial; it points to the subscriber's home
        network and is used by the GMSC to find the HLR.
  \item \textbf{IMSI} -- the subscription identity used inside the mobile network (MAP
        addressing, authentication); never dialled.
  \item \textbf{MSRN} -- a temporary E.164 number allocated by the visited VLR for one
        incoming call; it routes the call through the fixed network to the serving MSC.
\end{itemize}
The HLR maps MSISDN to IMSI and stores which VLR serves the subscriber; the VLR maps the
MSRN back to the IMSI when the call arrives.

\subsection{Location updating}
Three cases trigger a location update \cite{ts23012}:
\begin{itemize}
  \item \textbf{Normal} -- the MS reads a new location area identity (LAI) on the broadcast
        channel. If the new LA belongs to another VLR, the new VLR fetches the subscriber
        data from the HLR with \msg{Update Location}; the HLR sends \msg{Cancel Location} to the
        old VLR.
  \item \textbf{Periodic} -- the MS re-registers when timer T3212 expires, so the network
        can purge subscribers that disappeared without detaching (battery removed, coverage lost).
  \item \textbf{IMSI attach} -- at power-on in the same LA. The mirror procedure,
        \textbf{IMSI detach}, tells the VLR that the MS is off, so incoming calls are not paged
        in vain.
\end{itemize}

\subsection{Roaming}
A roaming subscriber registers in a visited network (VPLMN). The visited VLR addresses
the home HLR (HPLMN) over the international SS7 network, using global titles derived from
the IMSI, and receives the subscription data and authentication triplets. Roaming
agreements define charging and which services are allowed. Two consequences are worth
noting:
\begin{itemize}
  \item \textbf{Authentication stays with the home network}: only the home AuC knows $K_i$.
  \item \textbf{Tromboning}: a call to a roamer is always routed first to the home GMSC
        (the MSISDN belongs to the home network) and then internationally to the visited
        MSC. A local caller in the visited country therefore pays for two international
        legs, unless optimal routing is implemented.
\end{itemize}

\subsection{Call release}
Either party can release. The mobile sends \msg{Disconnect}; the MSC replies
\msg{Release} and the MS confirms \msg{Release Complete}. Towards the fixed network the MSC
sends ISUP \msg{REL} and receives \msg{RLC}, which frees the trunk. Finally the MSC orders
the BSS to release the radio channel (BSSMAP \msg{Clear Command}). Billing records are
closed when the circuit is released, which is why circuit-switched charging is naturally
time-based.

% ---------------------------------------------------------------------
\section{Mathematical models}

\subsection{The loss system and the Erlang-B formula}\label{sec:erlangb}
\paragraph{Engineering problem} A cell has $m$ traffic channels. Calls that arrive when all
are busy are \emph{blocked} (lost); they do not wait. How large must $m$ be so that no more
than, say, 2\% of attempts are blocked?

\paragraph{System model and assumptions} (i) Call arrivals form a Poisson process of rate
$\lambda$ from a large population; (ii) holding times are independent and exponentially
distributed with mean $h = 1/\mu$; (iii) blocked calls are cleared (no retry);
(iv) all $m$ channels are fully accessible. This is the $M/M/m/m$ queue.

\paragraph{Derivation} Let $p_k$ be the steady-state probability that $k$ channels are busy.
In state $k$ the system moves up at rate $\lambda$ and down at rate $k\mu$. Balancing the
flow between neighbouring states gives
\begin{equation}\label{eq:balance}
  \lambda\,p_{k-1} = k\mu\,p_k \quad\Longrightarrow\quad
  p_k = \frac{A}{k}\,p_{k-1} = \frac{A^k}{k!}\,p_0, \qquad A = \frac{\lambda}{\mu} = \lambda h .
\end{equation}
Normalising $\sum_{k=0}^{m}p_k = 1$ yields $p_0 = \bigl(\sum_{i=0}^{m}A^i/i!\bigr)^{-1}$.
Because Poisson arrivals see time averages (the PASTA property), the probability that an
arriving call is blocked equals the probability that all $m$ channels are busy:
\begin{equation}\label{eq:erlangb}
  P_{\mathrm{B}} = B(A,m) = \frac{A^m/m!}{\displaystyle\sum_{i=0}^{m}\frac{A^i}{i!}} .
\end{equation}
\paragraph{Symbols and units} $A$: offered traffic [E]; $m$: number of channels; $\lambda$:
arrival rate [calls/s]; $h$: mean holding time [s]; $P_{\mathrm{B}}$: blocking probability
(dimensionless). The carried traffic is $A_{\mathrm{c}} = A(1-P_{\mathrm{B}})$ and the
utilisation of each channel is $\rho = A_{\mathrm{c}}/m$.

\paragraph{Numerically stable form} Direct evaluation of \cref{eq:erlangb} overflows for
large $m$. Dividing numerator and denominator by $A^m/m!$ gives the recursion
\begin{equation}\label{eq:erlangrec}
  B(A,0) = 1, \qquad B(A,k) = \frac{A\,B(A,k-1)}{k + A\,B(A,k-1)}, \quad k = 1,\dots,m,
\end{equation}
used in \texttt{scripts/acnmath.py} to produce \cref{fig:erlang} and every number below.

\paragraph{Interpretation} \Cref{eq:erlangb} depends only on the \emph{mean} holding time
(it is insensitive to the holding-time distribution), which is why it worked well for
telephony for a century. It assumes, however, that blocked users go away; in reality many
retry, which increases the offered traffic and the true blocking (\cref{exr:retry}).

\begin{acnfigure}
  \acnfig{lecture-02/erlang-b}
  \caption{Erlang-B blocking probability versus offered traffic for $m = 8$, 30 and 94
  channels. For the same 2\% grade of service, 8 channels carry 3.6\,E (0.45\,E per channel)
  whereas 94 channels carry 82.2\,E (0.87\,E per channel).}
  \label{fig:erlang}
\end{acnfigure}

\subsection{Trunking efficiency}
\paragraph{Problem} Is it better to have one pool of $2m$ channels or two pools of $m$?
\paragraph{Result} For a fixed grade of service, large pools carry more traffic per channel.
This is the \emph{trunking gain} (or pooling gain) and follows from \cref{eq:erlangb}: the
relative fluctuation of the number of active calls, $\sqrt{A}/A$, shrinks as $A$ grows, so
less spare capacity is needed. \Cref{tab:trunking} quantifies it. It explains why GSM
operators avoid splitting channels between sectors unless interference requires it, and
why statistical multiplexing in packet networks is so attractive.

\begin{table}[htbp]
\centering\small
\caption{Traffic carried at 2\% blocking (computed with \cref{eq:erlangrec}).}\label{tab:trunking}
\begin{tabular}{@{}rrrr@{}}
\toprule
Channels $m$ & Offered traffic $A$ [E] & Carried traffic per channel [E] & Utilisation \\
\midrule
8  & 3.63  & 0.44 & 44\% \\
15 & 9.01  & 0.59 & 59\% \\
30 & 21.93 & 0.72 & 72\% \\
47 & 37.46 & 0.78 & 78\% \\
94 & 82.17 & 0.86 & 86\% \\
\bottomrule
\end{tabular}
\end{table}

\subsection{Waste of a reserved circuit}
If the voice activity factor is $\alpha$ (fraction of time a speaker is active), the
useful load of a reserved circuit is $\alpha$ times its capacity. With $\alpha\approx0.4$ in
each direction, roughly 60\% of every reserved circuit-second carries silence. Packet
switching recovers this capacity statistically, at the price of variable delay -- the
fundamental trade-off between the two switching paradigms, revisited in Lectures~3 and~5.

% ---------------------------------------------------------------------
\section{Network architecture and diagrams}
The SS7 network is a separate packet network overlaid on the circuit network. Every MSC,
GMSC and HLR is a signalling point; STPs route between them. Trunks between MSCs and the
PSTN carry voice, while MAP and ISUP messages travel on signalling links that may follow
completely different physical routes. \Cref{fig:mtc} shows the interplay for the hardest
case, an incoming call, and \cref{fig:moc} the outgoing case.

% ---------------------------------------------------------------------
\section{Step-by-step procedures}

\begin{procedurebox}{Mobile-originated call (MOC)}\label{proc:moc}
\begin{steps}
  \item The MS sends a channel request on the random access channel (RACH) and receives a
        signalling channel (SDCCH) on the access grant channel (AGCH).
  \item It sends \msg{CM Service Request}; the MSC/VLR authenticates the MS and starts
        ciphering (Lecture~1).
  \item The MS sends \msg{Setup} with the called number. The MSC checks the subscription
        (barring, prepaid credit) and answers \msg{Call Proceeding}.
  \item The MSC orders the BSS to assign a traffic channel (BSSMAP \msg{Assignment Request});
        the MS moves from SDCCH to a TCH.
  \item The MSC analyses the number and sends ISUP \msg{IAM} towards the destination; the
        remote exchange answers \msg{ACM} when the called phone rings, and the MSC sends
        \msg{Alerting} to the MS (ring-back tone).
  \item When the called party answers, ISUP \msg{ANM} arrives; the MSC sends \msg{Connect},
        the MS acknowledges, and charging starts.
\end{steps}
\end{procedurebox}

\begin{acnfigure}
  \resizebox{0.92\linewidth}{!}{\acnfig{lecture-02/moc-sequence}}
  \caption{Mobile-originated call (\cref{proc:moc}).}
  \label{fig:moc}
\end{acnfigure}

\begin{procedurebox}{Mobile-terminated call (MTC)}\label{proc:mtc}
\begin{steps}
  \item A caller in the PSTN dials the MSISDN. The PSTN routes the call (ISUP \msg{IAM}) to a
        GMSC of the home network.
  \item The GMSC sends MAP \msg{Send Routing Information} to the HLR.
  \item The HLR finds the serving VLR and requests a roaming number with MAP
        \msg{Provide Roaming Number}; the VLR allocates an MSRN and returns it.
  \item The HLR returns the MSRN to the GMSC, which routes the call with a new \msg{IAM}
        addressed to the MSRN -- the fixed network delivers it to the serving MSC.
  \item The MSC asks its VLR for the subscriber context; the VLR releases the MSRN and
        returns the location area; the MSC pages the MS in all cells of that LA.
  \item The MS answers the page, is authenticated, receives \msg{Setup} and a TCH; ringing
        and answer are signalled back with \msg{ACM} and \msg{ANM}.
\end{steps}
\end{procedurebox}

\begin{acnfigure}
  \resizebox{\linewidth}{!}{\acnfig{lecture-02/mtc-sequence}}
  \caption{Mobile-terminated call: HLR interrogation, MSRN allocation, routing and paging
  (\cref{proc:mtc}). Dashed: signalling; solid: the reserved circuit.}
  \label{fig:mtc}
\end{acnfigure}

% ---------------------------------------------------------------------
\section{Worked examples}

\begin{examplebox}{Dimensioning a cell}\label{ex:dim}
\textbf{Problem.} The cell of Lecture~1, Example~1, has 94 traffic channels. Subscribers
generate 25\,mE each during the busy hour. How many subscribers can the cell serve at a
2\% grade of service?

\textbf{Solution.} Solving $B(A,94) = 0.02$ with \cref{eq:erlangrec} gives
$A = 82.2$\,E. The number of subscribers is $N = A/a = 82.2/0.025 \approx 3290$.
Carried traffic $A_{\mathrm{c}} = 82.2\times0.98 = 80.5$\,E, i.e.\ 86\% channel utilisation.

\textbf{Interpretation.} If the same 94 channels were split into two sectors of 47, each
sector would carry 37.5\,E, 74.9\,E in total -- about 9\% fewer subscribers for the same
hardware (\cref{tab:trunking}).
\end{examplebox}

\begin{examplebox}{Erlang-B by hand}\label{ex:hand}
\textbf{Problem.} A small site has $m=3$ channels and $A=2$\,E. Compute $P_{\mathrm{B}}$.

\textbf{Solution.}
\begin{equation*}
  P_{\mathrm{B}} = \frac{2^3/3!}{1 + 2 + 2^2/2! + 2^3/3!}
  = \frac{1.333}{6.333} = 0.211 .
\end{equation*}
The recursion gives the same: $B_1 = 2/3$, $B_2 = (2\cdot\tfrac23)/(2+\tfrac43) = 0.4$,
$B_3 = 0.8/(3+0.8) = 0.211$.

\textbf{Interpretation.} One call in five is lost: a 3-channel site is useless at 2\,E.
Small pools need very low load for an acceptable grade of service.
\end{examplebox}

\begin{examplebox}{Busy-hour arrival rate}\label{ex:lambda}
\textbf{Problem.} Mean call duration is 90\,s. What arrival rate produces the 82.2\,E of
\cref{ex:dim}, and how many call set-ups per second must the MSC handle in that cell?

\textbf{Solution.} $\lambda = A/h = 82.2/90 = 0.913$ calls/s, i.e.\ about 3290 call attempts
per hour in one cell. Each attempt generates tens of signalling messages (RR, MM, CC,
BSSMAP and ISUP), so signalling load scales with $\lambda$, not with $A$.

\textbf{Interpretation.} Short calls (or SMS and location updates, which use SDCCH)
stress signalling capacity, which must be dimensioned separately (\cref{exr:sdcch}).
\end{examplebox}

% ---------------------------------------------------------------------
\section{Technology comparison}
\begin{table}[htbp]
\centering\small
\caption{Circuit switching versus packet switching.}\label{tab:csps}
\rowcolors{2}{acnmist}{white}
\begin{tabularx}{\linewidth}{@{}l Y Y@{}}
\toprule
\rowcolor{white} & Circuit switching (GSM voice) & Packet switching (GPRS and later) \\
\midrule
Resource allocation & Reserved at set-up for the whole call & Per packet or per scheduling interval \\
Delay               & Constant, low                          & Variable (queuing) \\
Loss mechanism      & Blocking at set-up                     & Packet loss or delay under load \\
Efficiency          & Low for bursty sources                 & High (statistical multiplexing) \\
Dimensioning model  & Erlang-B ($M/M/m/m$)                   & Queueing models ($M/M/1$, etc.) \\
Charging            & Time                                   & Volume or flat rate \\
Signalling          & SS7: ISUP + MAP                        & GTP, later Diameter and HTTP/2 \\
\bottomrule
\end{tabularx}
\end{table}

\paragraph{Why circuit switching was right for voice} Voice is a constant-rate, delay-sensitive
stream; reserving 13\,kbit/s guarantees quality without complex scheduling, and the Erlang
model allows precise dimensioning. \paragraph{Why it became a limitation} Data is bursty: a
web page needs a high rate for a second and nothing for a minute. Reserving a circuit for
that is wasteful, and a single 9.6\,kbit/s circuit is far too slow. Running two core
networks (circuit and packet) also doubled operational cost -- the motivation for the
all-IP core of LTE (Lecture~5), where voice itself becomes packets (VoLTE).

% ---------------------------------------------------------------------
\section{Key takeaways}
\begin{takeawaybox}
\begin{itemize}
  \item A GSM call reserves a radio timeslot, Abis/A channels, a switch path and trunks
        for its whole duration -- even during silence.
  \item SS7 separates signalling from voice: MTP transports, SCCP addresses (global titles),
        TCAP structures dialogues, MAP handles mobility, ISUP controls trunks.
  \item An incoming call needs the HLR (\msg{SRI}) and the visited VLR (\msg{PRN}) to obtain an
        MSRN before it can be routed and paged.
  \item Erlang-B, $B(A,m)$, dimensions loss systems; compute it with the recursion
        \cref{eq:erlangrec}.
  \item Larger pools are more efficient (trunking gain); splitting pools costs capacity.
  \item Circuit switching wastes the silent part of every call and cannot serve bursty
        data efficiently; this drove the move to packet switching.
\end{itemize}
\end{takeawaybox}

% ---------------------------------------------------------------------
\section{Exercises}

\begin{exercisebox}{1}{Signalling}\label{exr:ss7}
Why does ISUP run directly over MTP\,3, while MAP needs SCCP? Relate your answer to how
the destination of each message is known.
\end{exercisebox}

\begin{exercisebox}{1}{Numbers}\label{exr:msrn}
Why can the GMSC not route an incoming call using the IMSI? Why is the MSRN released as
soon as the call reaches the serving MSC?
\end{exercisebox}

\begin{exercisebox}{2}{Dimensioning}\label{exr:dim}
A cell has 30 traffic channels. (a) What traffic can it offer at 1\% and at 5\% blocking?
(b) With 20\,mE per subscriber, how many more subscribers are served at 5\% than at 1\%?
\end{exercisebox}

\begin{exercisebox}{2}{SDCCH capacity}\label{exr:sdcch}
Location updates, SMS and call set-ups each hold an SDCCH for about 3\,s. A cell sees 1.5
such transactions per second in the busy hour. How many SDCCH sub-channels are needed for
a blocking of at most 0.5\%?
\end{exercisebox}

\begin{exercisebox}{3}{Retries}\label{exr:retry}
Assume every blocked caller retries once, immediately and independently. Show that the
offered traffic becomes approximately $A' = A(1+P_{\mathrm{B}}')$ where
$P_{\mathrm{B}}' = B(A',m)$, solve this fixed-point equation numerically for $m = 30$,
$A = 21.93$\,E, and comment on the result.
\end{exercisebox}

\begin{exercisebox}{3}{Tromboning}\label{exr:trombone}
A subscriber of an Iranian operator roams in Turkey and is called by a Turkish fixed-line
user. Draw the circuit path and list which networks carry each leg. Propose a signalling
change that would avoid the international legs, and identify what information the
Turkish network would need.
\end{exercisebox}

\subsection*{Solutions}
\addcontentsline{toc}{subsection}{Solutions}

\begin{solutionbox}{\ref{exr:ss7}}
ISUP messages travel between adjacent exchanges along the circuit being set up; the next
exchange and its point code are known from the routing tables, so MTP\,3 routing is
enough. MAP messages are addressed to databases identified by subscriber identities (e.g.\
``the HLR of IMSI 432\ldots''), possibly in another country; SCCP global title translation
converts such an identity into a point code, hop by hop.
\end{solutionbox}

\begin{solutionbox}{\ref{exr:msrn}}
The PSTN routes by E.164 numbers; an IMSI is not a dialable number and is deliberately not
exposed outside the mobile network. The MSRN is a scarce temporary number taken from the
visited MSC's pool; once the call reaches that MSC the IMSI is retrieved from the VLR and
the MSRN can be reused for another call.
\end{solutionbox}

\begin{solutionbox}{\ref{exr:dim}}
(a) $B(A,30)=0.01 \Rightarrow A = 20.3$\,E; $B(A,30)=0.05 \Rightarrow A = 24.8$\,E.
(b) Subscribers: $20.3/0.02 \approx 1017$ versus $24.8/0.02 \approx 1240$, i.e.\ about 220
more (22\%) at the price of five times more blocking.
\end{solutionbox}

\begin{solutionbox}{\ref{exr:sdcch}}
$A = \lambda h = 1.5\times3 = 4.5$\,E. Iterating \cref{eq:erlangrec}: $B(4.5,10)=0.0105$,
$B(4.5,11)=0.0043$. Eleven sub-channels are needed; in practice operators configure two
SDCCH/8 channels (16 sub-channels), which also leaves room for peaks.
\end{solutionbox}

\begin{solutionbox}{\ref{exr:retry}}
A fraction $P_{\mathrm{B}}'$ of the $\lambda$ first attempts generate a second attempt, so
the total attempt rate is $\lambda(1+P_{\mathrm{B}}')$ and $A' = A(1+P_{\mathrm{B}}')$ (treating
retries as Poisson, which is an approximation). Iterating from $P_{\mathrm{B}}'=0.02$:
$A' = 22.37$, $B = 0.0236$; $A' = 22.45$, $B = 0.0243$; converging to
$A' \approx 22.47$\,E, $P_{\mathrm{B}}' \approx 0.0244$. The blocking seen by an attempt grows by
about 22\%; the fraction of \emph{users} who finally fail is $P_{\mathrm{B}}'^{\,2}\approx 0.06\%$,
but the network carries more signalling. Retries are a positive feedback that can
destabilise an overloaded network.
\end{solutionbox}

\begin{solutionbox}{\ref{exr:trombone}}
Path: Turkish caller $\to$ Turkish PSTN $\to$ international leg to the home GMSC in Iran
(the MSISDN is Iranian) $\to$ HLR interrogation returns a Turkish MSRN $\to$ second
international leg back to the visited MSC in Turkey $\to$ the roamer. Optimal routing
lets the Turkish gateway query the home HLR directly (MAP SRI towards the HPLMN) and route
to the MSRN inside Turkey; it requires the visited network to recognise the number as
belonging to a roaming subscriber and an agreement allowing the HLR query, and it changes
charging, which is why it was rarely deployed.
\end{solutionbox}

% ---------------------------------------------------------------------
\section{Further reading}
SS7 is introduced in \cite{ituq700}; MAP is specified in \cite{ts29002}, location
management in \cite{ts23012} and call handling in \cite{ts23018}. Layer-3 call control and
mobility messages are in \cite{ts24008}. Kleinrock \cite{kleinrock1975queueing} derives
the loss formula rigorously, and Rappaport \cite{rappaport2002wireless} applies it to
cellular dimensioning. Eberspächer et al.\ \cite{eberspaecher2009gsm} describe GSM call
flows in detail.

\printbibliography[heading=subbibliography,title={\ACNstrRefs}]
\end{document}
